\documentclass{situationist-preprint}

\title{Moral Luck in an Impeccable World}
\msid{REL-14-b}
\msfamily{REL-14}
\msdiscipline{Philosophy of Religion}
\mstype{Consequence}
\mskeywords{moral luck; impeccability; heaven; sourcehood; constitutive luck; circumstantial luck}

\begin{document}
\maketitle

\begin{abstract}
REL-14-a argues that the blessed are free as sources though they cannot sin, and that their character was formed by a history God did not settle. That history was shaped by circumstantial and constitutive luck: the blessed met the temptations they met and began with the dispositions they had. We derive two consequences. First, in the impeccable state circumstantial luck loses its moral bite, since no situation the blessed can be in leads them to wrong. Second, the constitutive luck that formed their character does not undermine their sourcehood or their praiseworthiness, by the distinction ETH-13-a draws between degree of culpability and scope of responsibility. Luck in the history is the price of the independence that Corollary 4.1 of REL-14-a requires.
\end{abstract}

\section{Two kinds of luck in heaven}

\citet{Nagel1979} distinguished resultant luck, in how one's acts turn out; circumstantial luck, in the situations one faces; and constitutive luck, in the temperament and dispositions one has. ETH-13 concludes that moral luck is real and does not undermine responsibility. ETH-13-a explains why: an agent's culpability, and by parity her praiseworthiness, depends only on the quality of will she exercised in acts under her control, while the scope of her responsibility extends to outcomes, circumstances and constitution that she did not control.

The blessed of REL-14-a are impeccable (Lemma 3.1) and free as sources but not as having alternatives to doing good (Theorem 3.3). By Corollary 4.1 of REL-14-a, their state is reachable only through a history God did not settle. Such a history is full of luck. One person met a grave temptation at the age at which she could least resist it; another never met it. One began with a quick temper, another with an even disposition. The question is whether this luck spoils either the freedom or the merit of the blessed.

\begin{proposition}\label{p:circ}
In the impeccable state, circumstantial luck makes no difference to whether the blessed do wrong.
\end{proposition}

\begin{proof}
By Lemma 3.1 of REL-14-a, there is no situation the blessed can be in in which they do wrong. Which situations they meet therefore makes no difference to whether they do wrong.
\end{proof}

Proposition~\ref{p:circ} marks a real difference between heaven and earth. On earth, two agents of equal quality of will may differ in what they do because one is tested and the other is not, and Williams's Gauguin shows how much can then turn on circumstance \citep{Williams1981}. In the impeccable state the test conditions are, in the relevant respect, uniform: every situation is one in which the good is apparent for what it is (Assumption 2.4 of REL-14-a). Circumstantial luck survives only as luck in which goods the blessed are given occasion to do, which bears on the scope of their goodness and not on its quality.

\begin{proposition}\label{p:const}
Constitutive luck in the formation of the blessed's character undermines neither their sourcehood nor their praiseworthiness for the acts by which that character was formed.
\end{proposition}

\begin{proof}
Sourcehood (Definition 2.1 of REL-14-a) requires that an action issue from a reasons-responsive mechanism that is the agent's own; it does not require that the agent controlled the origin of the dispositions on which the mechanism operates, any more than earthly responsibility does \citep[ch.~8]{FischerRavizza1998}. Praiseworthiness, by ETH-13-a, depends on the quality of will exercised in acts under the agent's control. The acts by which the blessed formed their characters were such acts, whatever dispositions they started from. Constitutive luck affects the scope of what they are responsible for, including how far each had to travel, and not the degree of their merit in the acts they performed.
\end{proof}

\section{Why the luck is not a defect}

\begin{objection}
If luck shaped the history, the difference between the blessed and the lost is partly a matter of luck, and a just God would not let eternal outcomes depend on it.
\end{objection}

\begin{reply}
The objection would have force against a God who sends some to endless punishment; REL-20 concludes that a just God leaves hell unoccupied. What remains is that the paths to the impeccable state differ in length and difficulty. That is a difference in scope, and ETH-13-a shows that differences in scope due to luck are consistent with responsibility. Nor could God remove the luck without removing the good at stake. A history whose circumstances and constitutions God fixed so as to guarantee each outcome would be one in which God settled the creatures' goodness, and by Corollary 4.1 of REL-14-a the independence of that goodness would then be lost \citep[cf.][]{PawlTimpe2009}.
\end{reply}

\section{Conclusion}

Moral luck is real in the history of the blessed and idle in their final state. Circumstantial luck can no longer lead them to wrong, and the constitutive luck that shaped their characters leaves their sourcehood and their merit intact. The luck in their history is not a flaw in the design of heaven. It is what a history God did not settle looks like from inside.

\begin{thebibliography}{99}\small
\bibitem[Fischer and Ravizza(1998)]{FischerRavizza1998} Fischer, J.~M. and Ravizza, M. (1998). \emph{Responsibility and Control: A Theory of Moral Responsibility}. Cambridge University Press.
\bibitem[Nagel(1979)]{Nagel1979} Nagel, T. (1979). Moral luck. In \emph{Mortal Questions}. Cambridge University Press.
\bibitem[Pawl and Timpe(2009)]{PawlTimpe2009} Pawl, T. and Timpe, K. (2009). Incompatibilism, sin, and free will in heaven. \emph{Faith and Philosophy} 26(4): 398--419.
\bibitem[Williams(1981)]{Williams1981} Williams, B. (1981). Moral luck. In \emph{Moral Luck: Philosophical Papers 1973--1980}. Cambridge University Press.
\end{thebibliography}

\end{document}
